\section{Consequences and scope}
\label{sec:scope}

The theorem describes two distinct compact invariant sets and their
associated measures. Keeping them separate prevents a false
unique-ergodicity assertion on a space containing finite cycles.

\begin{table}[ht]
\centering
\small
\begin{tabular}{@{}>{\raggedright\arraybackslash}p{0.20\textwidth}
 >{\raggedright\arraybackslash}p{0.34\textwidth}
 >{\raggedright\arraybackslash}p{0.36\textwidth}@{}}
\toprule
Object & Established structure & Role of the estimate \\
\midrule
$C=\cl_K\bigcup_e\Pi_e$
& Compact classical invariant set containing all old cycles
& Uniform high-level contacts give finite nets at every scale \\
\addlinespace
$\A=\lim_e\Pi_e$
& Hausdorff limit, disjoint from every old cycle, conjugate to $\Zp$
& Lower contacts give convergence; the finite average gives separation \\
\addlinespace
$\mu_e$ and $\mu$
& Full-sequence convergence to the unique invariant probability on $\A$
& Aligned orbit couplings give $W_\infty(\mu_d,\mu)\le|s|^{c_d}$ \\
\bottomrule
\end{tabular}
\caption{Exact implications for the containing closure, the limiting
system, and the cycle measures. Unique ergodicity belongs to $\A$,
not to $C$. All distances are classical absolute-value distances.}
\label{tab:objects}
\end{table}

For any real continuous function $\varphi$ on $C$, the proof also
gives the explicit modulus-of-continuity estimate
\[
 \left|\int\varphi\,d\mu_d-\int\varphi\,d\mu\right|
 \le \omega_\varphi\bigl(|s|^{c_d}\bigr).
\]
In particular a Lipschitz test function of constant $L$ has error
at most $L|s|^{c_d}$. The bound is for each fixed allowed multiplier;
neither its optimality nor uniformity near the boundary $|s|=1$ is
asserted.

The proof uses the source's optimal-cycle geometry, coprime Hensel
factorization, derivative products, and the elementary
characteristic-$p$ substitution identity. Its specific distribution
step is the all-anchor contact average with a bound diverging in
the anchor level. Compactness and the cyclic inverse-limit argument
then identify the resulting measure. No Galois irreducibility,
oriented character, or splitting-field nesting conclusion is needed
or obtained here. The adding-machine conjugacy is topological, not
an asserted analytic linearization.

The theorem does not address characteristic zero, characteristic two,
all normalized germs, or counting all periodic points of a polynomial.
The comparison with existing measure and ramification results is
bounded to the specified sources and does not certify worldwide
priority or the present status of every later or unpublished work.
The assertion proved here is the positive answer to the explicitly
identified version of the selected-cycle question.

\paragraph{Preparation and verification statement.}
This anonymous manuscript was prepared with AI assistance. The
underlying proof arguments received nonauthor review within the same
research team; this is internal verification, not external peer
review or journal acceptance. No numerical experiment is used as
evidence for the theorem. All substantive contact, compactness, and
aperiodicity arguments used in the article are included above.
